\scanpage{16}
\mspage{14}

3.\unsure{14|15}. Soit $\pi'$ un sous-groupe d'indice fini de $\pi$. Il correspond donc à une extension séparable finie $K'$ de $K$, d'où un normalisé $V'$ de $V$ dans $K'$, qui est un anneau de val. discrète hensélien. Si $k'$ est son corps résiduel, $k'/k$ est une extension finie, donc $\pi'_0\mathrel{\text{---}}\pi'$ \deleted{est} \inserted{induit} un isom de $\pi'_0$ \inserted{$=\operatorname{Gal}(\bar k'/k')$} sur un sous-groupe ouvert de $\pi_0=\operatorname{Gal}(\bar k/k)$. \deleted{D'autre part} \inserted{Il s'impose \unsure{de} vérifier que $\pi'_0\mathrel{\text{---}}\pi'$ est injectif, i.e. que} l'on a \deleted{on voit que}
\[
(3.\unsure{14}.1)\qquad I'=I\cap\pi'
\]
Notons que d'après la théorie générale appliquée à $V'$, on aura pour tt $\ell$ \inserted{nb premier} $\unsure{\ne}p$
\[
(\unsure{3.14.2|2.14.2})\qquad [I'/(I',I')]^{\deleted{\unsure{(\ell)}}(\ell)}\simeq\mathbb Z_\ell(1)
\]
(isom canonique). \deleted{D'autre part,} L'hom $I'_t\mathrel{\text{---}}I_t$ s'identifie, sur $I_t$, à la multiplication par $e$ \inserted{comme dessus}; \deleted{qui} lorsque $k'/k$ est séparable, \deleted{\unsure{p. ex. que}} i.e. l'extension résiduelle de $\widetilde V'/\widetilde V$ \deleted{séparable} \inserted{triviale}, \deleted{\unsure{donc}} p. ex. si $k$ est parfait, alors on peut interpréter $e$ en termes purement ``groupethéorique'' comme
\[
 e=[I:I']=\operatorname{card}I/I'.
\]
\inserted{Sinon, il faut diviser par le facteur inséparable $p^s$ de $k'/k$.}

\deleted{3.\unsure{16}. Il \unsure{reste à examiner}, lorsque \deleted{\unsure{$k'/k$}} \inserted{$I'$} \unsure{est invariant dans} $\pi$ \inserted{dans $I$}, $\unsure{\pi'/K|\pi'/I}$ \unsure{opère sur}}

3.15. Lorsque $\pi'$ est invariant dans $\pi$, \inserted{(i.e. $K'/K$ galoisienne)} il en est de même de \deleted{$\pi'_0$ dans $\pi_0$, puisque (on \unsure{sait} que $k'/k$ est \unsure{purement} galoisienne), donc \unsure{l'hom}} \inserted{$I'=I\cap\pi'$ dans $\pi$ (puisque $I$ est invariant)}, donc \deleted{$\unsure{\pi/\pi'_0}$} \deleted{\unsure{l'hom}} $\pi$ opère sur $I'/(I',I')$ \unsure{,} et par suite sur
\margin{15}
